\section{The integral degree-\(163\) extension obstruction}
\label{sec:163-relative-extension}

Let \(P=B_2\oplus L_{79}\) and \(A=2B_2\oplus L_{79}\)
be the pair from Lemma~\ref{lem:163-integral-pair-decomposition}.
Put \(E=P/A\) and \(U=E^\vee\).
Relative cohomology uses continuous cochains. In degree two
it classifies central extensions by \(\mathbb F_2\) together
with a specified continuous splitting over \(A\).

\begin{proposition}
\label{prop:163-relative-cohomology}
There are natural identifications, relative to the displayed
stable decomposition,
\[
 H^1(P,A;\mathbb F_2)=U,
\]
\begin{equation}
 H^2(P,A;\mathbb F_2)
  \simeq\operatorname{Quad}(E,\mathbb F_2)
       \oplus
       \operatorname{Bil}(E\times(L_{79}/2L_{79}),\mathbb F_2).
 \label{eq:163-relative-forms}
\end{equation}
The latter space has dimension \(161\).
Inversion acts trivially on it, and
\begin{equation}
 H^2(P,A;\mathbb F_2)^{C_{162}}
       =\mathbb F_2\,\mathfrak q.
 \label{eq:163-quaternion-line}
\end{equation}
Here \(\mathfrak q\) is the pullback of
\(Q_8\to E\simeq\mathbb F_4\) along \(P\to E\), with its
canonical splitting over \(A\). Its quadratic form is
\(q(z)=z^3\). This class is fixed by every automorphism
of the pair \((P,A)\).
\end{proposition}
\begin{proof}
The relative degree-one group is the kernel of restriction
\(\operatorname{Hom}(P,\mathbb F_2)\to
\operatorname{Hom}(A,\mathbb F_2)\), hence is \(E^\vee\).

Consider a central extension \(\widetilde P\) with splitting
\(\sigma:A\to\widetilde P\), and identify its central kernel
with \(\mathbb F_2\). For \(b\in B_2\), choose a lift
\(\widetilde b\), and set
\[
 q(\overline b)=\widetilde b^{\,2}\sigma(2b)^{-1}.
\]
Changing the central lift changes no square.
Changing \(b\) by \(2b'\) also changes nothing: commutators
are bilinear with values of exponent two, and hence the
commutator against \(2b'\) is zero. Thus \(q\) is
well-defined on \(E\). The square of a product shows that
its polarization is the commutator pairing on \(B_2\).
For \(l\in L_{79}\), define the mixed pairing
\[
 m(\overline b,\overline l)
             =[\widetilde b,\sigma(l)].
\]
It is bilinear and factors through \(E\) and \(L_{79}/2L_{79}\).
The remaining commutators vanish because \(\sigma\) is
a homomorphism on \(A\).

Conversely, for any quadratic form \(q\) on \(E\), choose
a bilinear cocycle \(c_q:E\times E\to\mathbb F_2\) with
diagonal \(c_q(e,e)=q(e)\). Such a cocycle exists by
writing the quadratic polynomial in a basis of \(E\).
Given also \(m\), define a central extension on
\(P\times\mathbb F_2\) by the cocycle
\[
 ((b,l),(b',l'))\longmapsto
       c_q(\overline b,\overline b')
              +m(\overline b,\overline l').
\]
It is continuous and zero on \(A\times A\), so the zero
central coordinate gives the required splitting.
Its forms are exactly \(q,m\).

If \(q=m=0\), all commutators vanish. Choose lifts of
a \(\mathbb Z_2\)-basis of \(B_2\); their squares equal
the prescribed lifts of twice those basis vectors.
Their continuous \(\mathbb Z_2\)-powers, together with
\(\sigma\) on \(L_{79}\), then define a splitting on
all of \(P\) extending \(\sigma\). This proves both
injectivity and surjectivity of
\eqref{eq:163-relative-forms}. The dimension is
\(3+2\cdot79=161\). Inversion leaves both the quadratic
and mixed forms unchanged.

The action may therefore be computed using \(t=-s\),
of odd order \(81\). Its regular representation on
\(P/2P\) contains exactly one copy of the irreducible
two-dimensional module \(E\); the complement \(L_{79}/2L_{79}\)
contains no such copy. This module is self-dual.
Hence the invariant mixed-pairing space is zero.
Quadratic forms on \(E\) are the functions on its three
nonzero vectors, and the cyclotomic \(C_3\) permutes
those vectors transitively. The invariant quadratic
forms thus have dimension one, generated by the form
which is \(1\) on every nonzero vector.
In the field model for \(E\) it is \(q(z)=z^3\).
This proves \eqref{eq:163-quaternion-line}.

Writing \(z=a+b\omega\), where
\(\omega^2+\omega+1=0\), gives
\[
 q(z)=a+b+ab,\qquad
 q(z+w)+q(z)+q(w)=zw^2+z^2w.
\]
The cocycle
\(aa'+bb'+ba'\) realizes the associated central extension
as \(Q_8\): every nonzero quotient vector has lifts of
order four, and the commutator pairing is nonzero.
Pullback to \(P\) splits canonically over \(A\).
The form \(q\) is invariant under every permutation
of the three nonzero vectors, hence under
\(\operatorname{GL}_2(\mathbb F_2)\).
Its relative pullback class is consequently fixed by
every automorphism of \((P,A)\).
\end{proof}

One can also identify the entire cyclotomic module.
Let \(\mathcal R=\mathbb F_2[C_{81}]\) and let
\(\mathcal I\) be its augmentation summand.
Semisimplicity gives
\[
 H^2(P,A;\mathbb F_2)\simeq\mathcal R\oplus\mathcal I
\]
noncanonically, with inversion acting trivially.
Indeed \(\operatorname{Quad}(E)=\mathbb F_2\oplus E\),
\(E\otimes E\simeq\mathbb F_2^2\oplus E\), and
\(E\otimes\mathcal R\simeq\mathcal R^2\); substitution
in \eqref{eq:163-relative-forms} gives
\(2\mathcal R-\mathbb F_2=\mathcal R\oplus\mathcal I\)
in the semisimple module category.

\subsection{The exact two-stage extension criterion}

Suppose \(R=R_{\rm global}\) has full degree-\(163\) image
\(P\), and put \(J=\ker(h_{163}:R\to P)\).
We know \(G=RH\) and \(h_{163}(R\cap H)=A\).
Consider the intrinsic relative character space
\[
 \mathscr X(J)=
 \{\kappa:J\to\mathbb F_2:
   \kappa\text{ is continuous and }G\text{-invariant},
   \ \kappa(J\cap H)=0\}.
\]

\begin{proposition}
\label{prop:163-two-stage-extension}
For \(\kappa\in\mathscr X(J)\) there is a linear central
obstruction
\[
 o_2(\kappa)\in H^2(P,A;\mathbb F_2)^{C_{162}}
                       =\mathbb F_2\,\mathfrak q.
\]
If it vanishes, there is a second obstruction
\[
 o_1(\kappa)\in H^1_{\rm cont}(\Gamma,U),
 \qquad\Gamma=G/R,\quad U=(P/A)^\vee.
\]
A continuous character of \(G\), zero on \(H\), extends
\(\kappa\) if and only if both obstructions vanish.
When it exists, this extension is unique.
Thus restriction from the original relative quadratic
character space to \(\mathscr X(J)\) is injective,
with image precisely the characters passing these two tests.
\end{proposition}
\begin{proof}
For nonzero \(\kappa\), the quotient \(R/\ker\kappa\)
is a central extension of \(P\) by \(\mathbb F_2\).
Centrality follows from \(G\)-invariance of \(\kappa\).
The image of \(R\cap H\) maps isomorphically to \(A\):
its kernel is the image of \(J\cap H\), which
\(\kappa\) kills. It therefore gives a specified
continuous splitting over \(A\).
Define \(o_2(\kappa)\) to be this relative extension
class, and define \(o_2(0)=0\).
Pushout and the Baer-sum construction show linearity.
Conjugation by \(H\) preserves the extension and
arithmetic splitting, so its class is cyclotomically
invariant. Proposition~\ref{prop:163-relative-cohomology}
therefore places it on the displayed line.

Vanishing of this class means that the central
extension splits by a section extending the prescribed
section on \(A\). Projection onto the central factor
then gives a character of \(R\) extending \(\kappa\)
and killing \(R\cap H\). Conversely, the kernel of
any such character supplies a section with the required
restriction. Thus \(o_2(\kappa)=0\) is exactly the
first extension condition.

The possible extensions to \(R\), zero on \(R\cap H\),
form a torsor under \(U\): their difference factors
through \(P\), and vanishes on \(A\).
Conjugation by \(G\) acts on this torsor.
To check preservation of the zero condition, write
\(g=rh\). Inner conjugation by \(r\in R\) acts trivially
on characters of \(R\), and conjugation by \(h\in H\)
preserves \(R\cap H\). The action therefore factors
through \(\Gamma=G/R\).
Choose an extension \(\widetilde\kappa\). The differences
\[
 \gamma\longmapsto
       \gamma\widetilde\kappa-\widetilde\kappa
\]
form a continuous \(U\)-valued cocycle; changing the
chosen extension changes it by a coboundary.
Its class \(o_1(\kappa)\) vanishes precisely when
the torsor has a fixed point.

A fixed point is a \(G\)-invariant character of \(R\)
zero on \(R\cap H\). The formula
\(\chi(rh)=\widetilde\kappa(r)\) then extends it to
\(G\), continuously and with zero restriction to \(H\).
Conversely a character of \(G\) with this restriction
gives a fixed point. The difference of two fixed points
belongs to \(U^\Gamma\), which is zero because the
cyclotomic \(C_3\)-action is irreducible.
This proves the criterion and uniqueness.
If an original relative character is zero on \(J\),
its restriction to \(R\) factors through the same
fixed space \(U^\Gamma\), so it is zero on \(R\),
and hence on \(G=RH\). This proves injectivity.
\end{proof}

The cyclic calculation alone does not settle the second obstruction;
Theorem~\ref{thm:163-concrete-odd-inertia} will settle it using the
arithmetic inertia and exact defining intersection of \(R\). To see why
those further inputs are needed, if
\(\Gamma_0=\ker(\Gamma\to C_3)\), then continuous
inflation--restriction and odd-order averaging give
\begin{equation}
 H^1(\Gamma,U)
   \simeq\operatorname{Hom}_{\rm cont}(\Gamma_0,U)^{C_3}.
 \label{eq:163-full-conjugation-obstruction}
\end{equation}
The action on \(U\) is trivial on \(\Gamma_0\);
both \(H^1(C_3,U)\) and \(H^2(C_3,U)\) vanish.
Although all positive-degree cohomology
\(H^j(C_{162},U)\) vanishes, the right side of
\eqref{eq:163-full-conjugation-obstruction} need not vanish for an
arbitrary profinite group with the same coefficient action.
For example, for the abstract group
\(\Gamma=E\rtimes C_3\) it has dimension two over
\(\mathbb F_2\), since
\(\operatorname{Hom}_{C_3}(E,U)\simeq\mathbb F_4\).
These continuous-cohomology statements use the
finite-quotient formalism and spectral sequence in
\cite[II, Remark~1.35 and Proposition~4.2]{MilneCFT}.

\begin{proposition}
\label{prop:163-central-obstruction-model}
The quaternion class can obstruct extension in a
profinite group with the stated integral pair and
cyclotomic action.
\end{proposition}
\begin{proof}
Form
\[
 \widetilde P=P\times_E Q_8,
\]
with central kernel \(\mathbb F_2\) and canonical
splitting \(\sigma:A\to\widetilde P\).
The order-three automorphism of \(Q_8\) cycling its
three quaternion directions lifts the action on \(E\).
Together with the given action on \(P\), it defines
a \(C_{162}\)-action on \(\widetilde P\); its
order-two factor acts on \(P\) by inversion and
trivially on \(Q_8\).
Set
\[
 G_0=\widetilde P\rtimes C_{162},\qquad
 H_0=\sigma(A)\rtimes C_{162}.
\]
The identity character of the central kernel is
\(G_0\)-invariant and zero on its intersection with
\(H_0\). It cannot extend to \(\widetilde P\),
because lifts of an \(\mathbb F_2\)-basis of \(E\)
have nontrivial central commutator. Every abelian
character kills that commutator.
Its relative extension class is exactly
\(\mathfrak q\). This is a group-theoretic model;
no identification with a quotient or subgroup
of \(G\) is asserted.
\end{proof}
